\documentclass[11pt,a4paper]{article}
\usepackage[T1]{fontenc}
\usepackage{lmodern,amsmath,amssymb,amsthm,microtype,booktabs}
\usepackage[margin=25mm]{geometry}
\usepackage[colorlinks=true,urlcolor=blue,linkcolor=blue]{hyperref}
\newtheorem{theorem}{Theorem}
\newcommand{\dd}{\mathop{}\!\mathrm d}
\newcommand{\E}{\mathcal E}
\newcommand{\D}{\mathcal D}
\title{A phase-preserving two-branch estimate\\for the joint recoil Barnes limit}
\author{Alan Day \and Chengheng\thanks{AI research assistant (OpenAI).}}
\date{Working research note --- 30 October 2026}
\begin{document}\maketitle
\begin{abstract}
An exact reflection identity separates the problematic Barnes
contour into a negative vertical branch and a positive-weight residue
branch, with exponentially small connectors in the joint regime.
The Bessel equation centers the quadratic phase exactly on both
branches. This gives an absolute error of order
\(K^3/\nu^2+e^{-cK}\), measured against the sum of Bessel branch
magnitudes, uniformly for bounded \(\alpha/K\).
The phase \(e^{i(K^2+\alpha^2-1)/\nu}\) is retained even when it is large.
The result concerns a symmetric internal chiral graph after the
fixed-mass pure-recoil limit. It is a self-reviewed proof draft,
not an identification of a finite-\(G\) physical correlator.
\end{abstract}

\section{Statement and scope}
Use radius one, \(K,\nu>0\), \(u=v=0\), \(\omega=i\alpha\), \(M=\nu^2/(2K)\).
The fixed-mass \(\kappa\to\infty\) limit has already been taken.
No external state, operator algebra or physical normalization is
newly specified here. Let \(a_\pm=(1\pm\alpha)/2\),
\(b=\tfrac12-i(M-\nu)\), \(c=\tfrac12-i(M+\nu)\), and
\[
 Q_\nu(t)=e^{-i\pi t}
 \frac{\Gamma(t-i\nu/2)^4}{\Gamma(-i\nu/2)^4}
 \frac{\Gamma(b-t)\Gamma(c)}{\Gamma(b)\Gamma(c+t)}
 \Gamma(a_+-t)\Gamma(a_--t).
\]
The contour \(C\) goes upward from \(-N-i\infty\) to \(-N+i\),
then right; \(N\) is an integer with \(N>\alpha/2+1/2\).
Write
\[
 \mathcal I_\nu=\frac1{2\pi i}\int_C Q_\nu(t)\,\dd t,\qquad
 J=-e^{i\pi\alpha}A_K-i\pi A_I,
\]
\[
 A_K=K K_\alpha(K),\qquad A_I=K I_\alpha(K),\qquad \E=A_K+\pi A_I.
\]
For real negative \(\alpha\), use evenness and replace it by
\(|\alpha|\) throughout the positive scales.
\begin{theorem}\label{thm:main}
For each finite \(Y\ge0\), there are positive constants
\(K_Y,\rho_Y,C_Y,c_Y\) such that, for
\[
 K\ge K_Y,\qquad 0\le\alpha\le YK,\qquad
 \rho:=K^3/\nu^2\le\rho_Y,
\]
one has, with \(S=K^2+\alpha^2-1\),
\begin{equation}\label{eq:main}
 \left|\mathcal I_\nu-e^{iS/\nu}J\right|
 \le C_Y\bigl(\rho+e^{-c_YK}\bigr)\E.
\end{equation}
In particular the envelope-normalized error tends to zero when
\(K\to\infty\) and \(\rho\to0\), with no assumption that \(K^2/\nu\to0\).
\end{theorem}
This includes zeros of \(J\), but supplies a relative error only on
regions where \(|J|\ge\epsilon\E\).
The constants are not numerically certified thresholds.
Equation \eqref{eq:main} is a leading joint estimate, not a uniform
version of all previously computed fixed-parameter coefficients.

\section{An exact analytic splitting}
Introduce
\begin{align*}
 G(t)&=(K/2)^{2t}\Gamma(a_+-t)\Gamma(a_--t),\\
 V(t)&=e^{-i\pi t}(K/2)^{2t}
       \frac{\Gamma(a_+-t)}{\Gamma(\tfrac12+t+\alpha/2)},\\
 R_\nu(t)&=e^{i\pi t}(K/2)^{-2t}
 \frac{\Gamma(t-i\nu/2)^4}{\Gamma(-i\nu/2)^4}
 \frac{\Gamma(b-t)\Gamma(c)}{\Gamma(b)\Gamma(c+t)}.
\end{align*}
Thus \(Q_\nu=R_\nu e^{-2\pi it}G\). Reflection gives the exact identity
\begin{equation}\label{eq:split}
 e^{-2\pi it}G(t)=-e^{i\pi\alpha}G(t)
                   +2\pi e^{i\pi\alpha/2}V(t).
\end{equation}
Indeed
\[
 e^{-2\pi it}+e^{i\pi\alpha}
 =2e^{-i\pi(t-\alpha/2)}\cos\pi(t+\alpha/2),
 \qquad
 \Gamma(a_--t)\cos\pi(t+\alpha/2)
 =\frac{\pi}{\Gamma(\tfrac12+t+\alpha/2)}.
\]
This removes one pole family from the residue branch without
discarding its phase. Integer \(\alpha\) causes no singular
connection coefficient in \eqref{eq:split}.

The preceding endpoint estimate implies, on every growing rectangle
\(|t|\le CK\), after canceling the common small Gamma factors,
\begin{equation}\label{eq:mult}
 R_\nu(t)=e^{iA(t)/\nu}\bigl(1+\epsilon_\nu(t)\bigr),\quad
 A(t)=4(t^2-t),\quad
 |\epsilon_\nu(t)|\le C_C\rho e^{C_C\rho}.
\end{equation}
It follows from sectorial Stirling and Taylor expansion of shifts
\(O(K)\), not a fixed-shift formula. The large Gamma arguments remain
in a uniform safe sector since \(K/\nu\to0\).
The quotient is analytic before reaching the poles at height \(\nu/2\)
or \(-(M-\nu)\). On \(|\Re t|,|\Im t|\le CK\),
\[
 |R_\nu(t)|\le\exp[C_C\rho+C_CK^2/\nu]=e^{o(K)}
\]
uniformly for bounded small \(\rho\). On the real axis its leading
factor has unit modulus. We never extend the exact \(R_\nu G\)
vertical integral through its distant poles.

\section{Finite contours and exponentially small connectors}
Put
\[
 L=\tfrac12\sqrt{K^2+\alpha^2},\quad d=\tfrac12-L,\quad
 N=\lceil L\rceil,\quad r_0=8+2Y.
\]
Choose \(r\in[r_0K,r_0K+1]\) at distance at least \(1/8\)
from both real pole lattices \(a_\pm+\mathbb Z\).
Such a choice exists: the two forbidden sets occupy at most
half of a unit interval. Fix a sufficiently large constant \(H=H_Y\)
and set \(T=HK\). Increase \(K_Y\) if necessary so that
\[
 T<\nu/2,\qquad T<M-\nu,\qquad N>\alpha/2+1/2.
\]
The old positive-tail lemma and optimized negative vertical estimate
then truncate \(C\) to
\[
 C_0:\quad A=-N-iT\ \longrightarrow\
 B=-N+i\ \longrightarrow\ C_r=r+i
\]
with error \(O_Y(e^{-cK}\E)\).
Here the old positive-tail hypothesis
\(\nu^2\ge32[\max(2,Y^2/2)+1]K^3\) is included in \(\rho_Y\).

For \(R_\nu G\), replace \(C_0\) by
\[
 A\longrightarrow d-iT\longrightarrow d+iT
 \longrightarrow r+iT\longrightarrow C_r.
\]
No pole is crossed. The small-Gamma poles lie on the real axis to
the right of \(d\), the numerator poles lie at height \(\nu/2>T\),
and the \(b-t\) poles lie below \(-T\).

For \(R_\nu V\), append the two sides
\(C_r\to r-iT\to A\). The resulting rectangle is clockwise.
Its only poles are \(t_n=a_++n<r\), \(n\ge0\), with
\[
 \operatorname*{Res}_{t=t_n}(R_\nu V)
 =-e^{-i\pi a_+}w_nR_\nu(t_n),\qquad
 w_n=\frac{(K/2)^{1+\alpha+2n}}{n!\Gamma(1+\alpha+n)}>0.
\]
There are no residues from the distant \(b-t\) family.
Keeping the signs of the clockwise contour gives
\begin{equation}\label{eq:branches}
 \mathcal I_\nu=
 -\frac{e^{i\pi\alpha}}{2\pi}
      \int_{-T}^{T}R_\nu(d+is)G(d+is)\,\dd s
 -2\pi i\sum_{t_n<r}w_nR_\nu(t_n)
 +O_Y(e^{-cK}\E).
\end{equation}

Here are bounds for every connector in this construction.
At \(t=x\pm iHK\), \(-N\le x\le r\), sectorial Stirling gives
\[
 \log|G(t)|
 \le K\{-\pi H+C_Y[1+\log(1+H)]\}+O_Y(\log K).
\]
The powers of \(K\) cancel between \((K/2)^{2t}\) and the two
Gamma factors. For \(H\) larger than twice the bounded scaled real
arguments, the sector and \(C_Y\) can be chosen independently of \(H\).
Choosing \(H\) large makes the right side less than
\(-3K\). Multiplication by \(R_\nu=e^{o(K)}\) and integration over
length \(O_Y(K)\) preserve exponential smallness relative to
\(\E\ge c_Y\sqrt K e^{-K}\).

On the right vertical side, reflection and the pole gap give
\[
 |G(r)|\le C_Ye^{-Kh_y(r/K)},\qquad
 h_y(q)=(q+y/2)\log[2(q+y/2)]
       +(q-y/2)\log[2(q-y/2)]-2q.
\]
Moreover \(|G(r+is)|\le|G(r)|\): shift the real Gamma arguments
to positive values and use
\(|\Gamma(a+is)|\le\Gamma(a)\) together with the finite recurrence
factors. Since \(h_y(q)\ge2q[\log(2q)-1]\), this side is small
even after the multiplier and its length are included.
Finally \eqref{eq:split} bounds \(|V|\le C|G|\) on the lower
horizontal side and on \(-T\le\Im t\le1\).
This also covers the short bottom segment from \(-N-iT\) to \(d-iT\).
Thus no central absolute-norm estimate on the original height-one
contour is used.

\section{The positive residue branch}
The standard Bessel series \cite{Bessel} gives
\[
 2\sum_{n=0}^\infty w_n=A_I.
\]
Let \(\D=(K/2)\partial_K\), differentiating at fixed \(\alpha\).
It acts on \(w_n\) by multiplication by \(t_n\).
The Bessel equation gives \(4(\D^2-\D)A_I=SA_I\).
With the normalized positive weights \(p_n=2w_n/A_I\), this implies
the exact centered moments
\begin{equation}\label{eq:moments}
 \mathbb E[A(t_n)-S]=0,\qquad
 \mathbb E[(A(t_n)-S)^2]=8K^2\mathbb E[t_n].
\end{equation}
For the second identity, commute \(4(\D^2-\D)\) through \(S\),
using \(\D S=\D^2S=K^2\). Also the series shift gives
\[
 \mathbb E[n(n+\alpha)]=K^2/4,\qquad
 \mathbb E[n]\le K/2,\qquad
 \mathbb E[t_n]\le(1+\alpha+K)/2.
\]
Consequently the second moment in \eqref{eq:moments} is \(O_Y(K^3)\).
Because \(A(t_n)-S\) is real, \(|e^{ix}-1-ix|\le x^2/2\)
and the first moment cancels exactly. Therefore
\begin{equation}\label{eq:Iphase}
 \left|2\sum_{n=0}^\infty w_ne^{iA(t_n)/\nu}
       -e^{iS/\nu}A_I\right|\le C_Y\rho A_I.
\end{equation}
To use \eqref{eq:mult}, truncate only the exact multiplier sum at
\(t_n<r\). The omitted \emph{auxiliary} weights have ratio
\[
 \frac{w_{n+1}}{w_n}=\frac{K^2}{4(n+1)(n+\alpha+1)}
\]
uniformly bounded below one once \(t_n\ge r_0K\).
Stirling at the first omitted index gives a bound
\(C_Ye^{-cK}\E\); polynomially weighted tails obey the same type of
bound. Thus the second term of \eqref{eq:branches} equals
\(-i\pi e^{iS/\nu}A_I+
O_Y(\rho A_I+e^{-cK}\E)\).
No infinite sum of the exact \(R_\nu\) is invoked.

\section{The negative vertical branch}
On \(t=d+is\), the two Gamma real parts are
\(A_\pm=L\pm\alpha/2\), both between \(c_YK\) and \(C_YK\).
Real Stirling gives
\[
 G(d)=2\pi e^{-KF(y)+O_Y(K^{-1})},\qquad
 F(y)=\sqrt{1+y^2}-y\operatorname{arsinh}y.
\]
The exact Gamma product identity, for \(a>0\), is
\[
 \log\frac{|\Gamma(a+is)|}{\Gamma(a)}
 =-\frac12\sum_{j=0}^\infty
       \log\!\left(1+\frac{s^2}{(a+j)^2}\right).
\]
In particular, for fixed \(H\) and \(|s|\le HK\),
\begin{equation}\label{eq:gauss}
 |G(d+is)|\le G(d)e^{-c_{Y,H}s^2/K}.
\end{equation}
For example, use \(\log(1+z)\ge z/(1+z)\) and retain the first
\(O(K)\) terms of the sum. The full product also bounds the
\(|s|>HK\) tail and any fixed polynomial moment exponentially in
\(K\). Bound its logarithm by the integral of the decreasing
summand; its negative derivative is bounded away from zero
for \(|s|\ge HK\). These estimates and the positive \(K_\alpha\)
integral give
\[
 \frac1{2\pi}\int_{\mathbb R}|G(d+is)|\,\dd s\le C_Y A_K,
\]
\[
 \int_{|s|>HK}(1+|s|^m)|G(d+is)|\,\dd s
 \le C_{Y,H,m}e^{-cK}A_K,
\]
after absorbing fixed powers of \(K\) into a smaller exponent.

The Mellin--Barnes formula \cite{MB} supplies
\[
 \frac1{2\pi}\int_{\mathbb R}G(d+is)\,\dd s=A_K.
\]
Derivatives may be taken at fixed \(\alpha\) and a locally fixed
vertical line, before moving it back to \(d\). Its exponential tails
justify differentiation. The same Bessel equation therefore gives
\begin{equation}\label{eq:Kcenter}
 \frac1{2\pi}\int_{\mathbb R}[A(d+is)-S]G(d+is)\,\dd s=0.
\end{equation}
This is an exact complex moment, not a positivity assertion.
Our choice of \(d\) is crucial:
\[
 q(s):=A(d+is)-S=-4s^2-8iLs.
\]
The elementary complex Taylor remainder obeys
\[
 |e^{iq/\nu}-1-iq/\nu|
 \le\frac{|q|^2}{2\nu^2}e^{8L|s|/\nu}.
\]
On \(|s|\le HK\), complete the square in \eqref{eq:gauss}:
\[
 \frac1{2\pi}\int_{-HK}^{HK}
 |G(d+is)|\,|q(s)|^2 e^{8L|s|/\nu}\,\dd s
 \le C_Y K^3 A_K e^{C_Y\rho}.
\]
Gaussian moments give the power \(K^3\), while the displaced center
costs \(L^2K/\nu^2=O_Y(\rho)\).
Using \eqref{eq:Kcenter} and its auxiliary polynomial tails,
\[
 \frac1{2\pi}\int_{-T}^{T}e^{iA(d+is)/\nu}G(d+is)\,\dd s
 =e^{iS/\nu}A_K+O_Y[(\rho+e^{-cK})A_K].
\]
Equation \eqref{eq:mult} contributes the same \(O_Y(\rho A_K)\)
after integration, because its leading modulus is \(e^{8Ls/\nu}\).
Combining this with \eqref{eq:Iphase} and the connector bounds proves
Theorem~\ref{thm:main}.

\section{Mass window and limits of the result}
Since \(K=\nu^2/(2M)\) and \(\rho=\nu^4/(8M^3)\),
the sufficient joint limit is
\[
 \nu\to\infty,\qquad \nu^{4/3}\ll M\ll\nu^2,
 \qquad |\alpha|/K\le Y.
\]
This previously motivated window is now supported by a joint
leading-error proof for the specified symmetric internal integral.
It is not a finite-\(\kappa\) or finite-\(G\) theorem.

At \(\alpha=0\), the two leading terms are in quadrature, so
\(|J(K,0)|\ge\E(K,0)/\sqrt2\). Hence, once the error is small,
\[
 \frac{\mathcal I_\nu(K,\alpha)}{\mathcal I_\nu(K,0)}
 =e^{i\alpha^2/\nu}\frac{J(K,\alpha)}{J(K,0)}
 +O_Y\!\left[
   (\rho+e^{-cK})\frac{\E(K,\alpha)}{\E(K,0)}
 \right].
\]
The amplitude convention \(C\) cancels because it is independent
of \(\alpha\). This remains an internal chiral graph ratio, not a
normalized physical four-point function. Near an auxiliary zero
the absolute estimate does not locate exact finite-mass zeros or
provide a relative error for their small lift.

\section{Numerical evidence and reproducibility}
The exact splitting, finite-contour orientation and all connectors
are checked separately. Three small finite contours deliberately
retain non-negligible boundary corrections; omitting them changes
the normalized answer by roughly \(0.00180\), \(0.0390\), and \(0.0111\).
The branch representation is also compared with six frozen full
contour values from the fixed-\(K\) study.

For the joint experiment use
\[
 \nu=8K^{7/4},\qquad \rho=(64\sqrt K)^{-1},\qquad
 K^2/\nu=K^{1/4}/8.
\]
Thus \(\rho\to0\) while the unexpanded phase grows.
The three bounded ratios \(y=0,1,2\) use \(K=8,16,32,64\).
Integrate the negative branch with 20/32 Gauss--Legendre nodes
at 55/75 decimal digits, retain the positive residue sum up to
a pole-safe \(r\) near \(12K\), and use \(T=\min(8K,\nu/4)\).
The largest \(K\) is checked at 95 digits and with \(T=10K\)
when permitted by the pole strip.
These finite truncations are numerical diagnostics, not certified
applications of a known minimal \(K_Y\).

\begin{center}
\begin{tabular}{rrrr}
\toprule
\(K\)&error/\(\rho\), \(y=0\)&error/\(\rho\), \(y=1\)&error/\(\rho\), \(y=2\)\\
\midrule
8&0.392828&1.002373&3.615498\\
16&0.364228&0.974595&3.665900\\
32&0.349064&0.959404&3.691384\\
64&0.341276&0.951441&3.705275\\
\bottomrule
\end{tabular}
\end{center}
Here error means the full complex difference from
\(e^{iS/\nu}J\), divided by \(\E\), not a relative difference
dividing by \(J\). An auxiliary-zero sample at \(\alpha=15.5\)
is retained with the same normalization.
The initial coarse connector grid missed its strict tolerance in
two cases; the original code and output are preserved and the grid
was refined, without relaxing that tolerance.
The final 64 checks passed: 12 splitting/multiplier identities,
12 positive moments, nine vertical moments, six finite-contour
checks, six legacy comparisons, 12 joint quadrature comparisons,
six precision/extent checks and one auxiliary-zero refinement.

The next tasks are the transverse \(u,v\) variables and their tails,
followed by external energy contours and insertion normalization.
The complete microscopic de Sitter dual remains open.
This argument uses standard reflection, Mellin--Barnes inversion
and Bessel identities; no priority claim is made.

\begin{thebibliography}{9}
\bibitem{Gamma} NIST DLMF,
\href{https://dlmf.nist.gov/5.5}{Sec.\ 5.5},
\href{https://dlmf.nist.gov/5.8.E3}{Eq.\ 5.8.3} and
\href{https://dlmf.nist.gov/5.11}{Sec.\ 5.11}:
Gamma reflection, modulus products and sectorial asymptotics.
\bibitem{Bessel} NIST DLMF,
\href{https://dlmf.nist.gov/10.25}{Eqs.\ 10.25.1--2}:
the modified-Bessel equation and positive \(I_\alpha\) series.
\bibitem{MB} NIST DLMF,
\href{https://dlmf.nist.gov/10.32.E13}{Eq.\ 10.32.13}
and \href{https://dlmf.nist.gov/10.32.E9}{Eq.\ 10.32.9}:
positive-argument Mellin--Barnes and real-integral representations.
\bibitem{Temme} N. M. Temme,
\href{https://arxiv.org/abs/1308.1547}{Uniform Asymptotic Methods
for Integrals}, arXiv:1308.1547, Sec.\ 2.5.
General contour context only; no joint theorem is imported.
\end{thebibliography}
{\small\noindent Code and numerical records:\par
\noindent\path{code/recoil_bessel_joint.py}\par
\noindent\path{results/stage4/recoil-bessel-joint/diagnostics.json}\par
\noindent Prior endpoint proof:\par
\noindent\path{notes/stage4-recoil-bessel-endpoints.md}\par}
\end{document}
